Helena
----------

Helena Carbajosa ha dado de qué hablar toda su vida. Su traducción de la Odisea, publicada a los 25 años, revolucionó el campo por su osadía\footnote{Le da un aíre a la traducción de un tal JM Tobal en nuestro universo paralelo}. A los 27 tiene un éxito similar con la Ilíada. A los treinta, su traducción de los poemas de Safo la confirman como la helenista más original en español y le valen, dos años más tarde, una cátedra en la U. de Ginebra, que acepta, abandonando su plaza en la Complutense y divorciándose a la vez de su marido, León Suárez, catedrático en la U. Autónoma de Madrid, y considerado el gran pope del Helenismo español. 

Suárez dirigió la tesis de Helena y facilitó tanto la publicación de sus traducciones como su difusión a nivel nacional e internacional. Veinte años mayor que ella, ha sido su protector y maestro, pero la relación entre ambos no se hizo pública ---a pesar de que las malas lenguas aseguran que empezó cuando ella cursaba el primer año de filología latina y griega--- hasta después de que Helena obtuviera un puesto de profesora titular en la Autónoma, a la tierna edad de 28 años.

 Durante los cuatro años que dura su matrimonio, Helena y León han sido la pareja de moda de las letras españolas. Suárez no sólo es un reputado filólogo, sino también un célebre poeta. La pareja es admirada y vilipendiada a partes iguales, sobre todo a partir de que Helena gana el premio Princesa de Asturias. Aunque León no forma parte del jurado el año que Helena ---con treinta primaveras--- se lleva el premio, ha sido su presidente hasta poco tiempo atrás y se sabe que su influencia es enorme. Sin embargo, nadie se atreve a cuestionarlos en público. León es poderoso e influyente y ha demostrado en numerosas ocasiones que puede hundir a un enemigo sin pestañear. Y el trabajo de Helena es verdaderamente brillante. 
 
 El divorcio y desplante de Helena causa un auténtico tsunami de cotorreos en los círculos ilustrados, de los que se hacen eco también los influencers de moda. Helena ha mantenido toda su carrera un perfil público muy alto. Físicamente es muy atractiva: pelirroja, ojos verdes, muy atlética, andrógina, élfica. No se prodiga en las redes sociales pero tampoco es tímida y sus opiniones tienden a ser muy iconoclastas. Entre ellas: su desprecio manifiesto por la mediocridad de la endogámica sociedad española ---sin que parezca importarle que en su propia carrera se la pueda acusar de haberse aprovechado de un marido influyente---; sus puntos de vista radicales como traductora ---defiende a capa y espada la necesidad de recuperar la <<esencia>> de los clásicos, de ahí sus traducciones en hexámetros del griego y su introducción constante de atrevidos neologismos--- y sus posiciones políticas, muy izquierdistas. Sin embargo, Helena defiende celosamente su vida privada, limitándose a compartir en las redes información pública y siempre <<oficial>> de su relación con León. Con el divorcio, arrecian los rumores de una vida secreta, que incluye aventuras más o menos turbias, involucrando relaciones bisexuales, aparentemente toleradas por su ya ex-marido.  
 
 Otros rumores aseguran que el divorcio no tiene nada que ver con su vida sentimental sino con la necesidad de librarse de la influencia de León y establecer su propia voz como poeta. Helena solo ha publicado traducciones, aunque los rumores insisten en que algunos poemas de extraordinaria calidad que han circulado por páginas literarias en Internet sin firma, son de ella. 
 
Por último, se especula mucho sobre sus razones para elegir Ginebra, en lugar de Oxford, que también le había hecho una oferta. Las razones parecen ser claramente personales. Helena creció en Ginebra, hija de funcionarios internacionales y no vivió en España hasta los dieciocho años, cuando decidió estudiar en Madrid. De hecho, tiene un ligero acento en español, apenas audible y, según sus críticos, postizo. 

Theopisti
-------------

Penélope Davros, <<Penny>>, tiene muchas cosas en común con Helena. Ha crecido también en Ginebra, de padre griego y madre estadounidense, propietarios de un famoso restaurante. También abandona la ciudad en su adolescencia, para estudiar en Harvard, donde se convierte en la pupila ---y según rumores muy informados--- también en la amante de Theodora Kallis, la famosa papiróloga norteamericana de origen cretense.  Theodora es la máxima autoridad mundial en la obra de Safo
 

Lucía Arrieta, 44 años, es helenista en Salamanca. Su traducción de Safo es la de referencia en español. A los veinticinco publicó un único libro de poemas, hoy descatalogado, y no ha vuelto a escribir. Traduce porque dejó de atreverse a decir nada con su propia voz.
Delia Mavros, 36 años, es hispano-griega y restauradora de papiros en un laboratorio privado de Ginebra. Es la falsificadora.
Su pasado, fijado de antemano: se formó con Theodora Kallis, una papiróloga cretense veinte años mayor, que fue su maestra y su amante. Theodora murió hace nueve años obsesionada con Safo. Entre sus papeles había un fragmento auténtico, pequeño e inédito, de origen turbio. Delia empezó a completarlo como un duelo: escribir a Theodora en la lengua que las unía. Simon Achterberg, marchante de antigüedades en Ginebra, descubrió lo que hacía y lo convirtió en negocio. Tiene sobre ella pruebas de la sustracción de papiro antiguo en blanco y la sostiene con esa amenaza.
Conrad Weller es un magnate estadounidense con una fundación y un museo en construcción. Necesita una gran pieza para la inauguración.
Graham Pell es papirólogo en Oxford, rival académico de Lucía. Es un escéptico profesional y a veces tiene razón.
Primer acto
La fundación Weller adquiere a través de Achterberg un papiro que completa un poema conocido de Safo, y contrata a Lucía para autentificarlo y traducirlo de cara a la inauguración del museo, dentro de cinco meses. Achterberg anuncia que hay más fragmentos de la misma procedencia. Ese es el plazo.
Lucía viaja a Ginebra y examina el papiro en el laboratorio donde trabaja Delia. El poema la deja sin aliento. El lector nunca lo lee: solo ve lo que le hace a ella. Insomnio, la necesidad de caminar y una frase que no consigue traducir. La autentificación material es impecable, porque el papiro y la tinta son antiguos. Lucía firma el informe preliminar.
Entre las dos surge una atracción inmediata, primero técnica y luego otra cosa. Delia conoce su traducción de memoria.
Segundo acto
Van llegando nuevos poemas, uno cada pocas semanas, cada uno más intenso que el anterior. Weller sube la apuesta y planea una edición mundial. Pell publica un artículo que pone en duda la procedencia, y la reputación de Lucía queda atada a la autenticidad de los poemas. Lucía y Delia se acercan durante las sesiones de trabajo, y la relación empieza en un viaje a Lesbos que Weller paga para un documental.
Lucía empieza a ver grietas que primero racionaliza. Una palabra con rasgos cretenses en lugar de eolios. Una planta, el díctamo, que no crece en Lesbos. Un giro sintáctico que no consigue documentar. Cada anomalía tiene una explicación posible, y ella elige esa explicación, como elige ciertas lecturas un enamorado. Lo que la convence no es la filología, sino que los poemas le hablan a ella.
En el punto medio, Pell le hace llegar, con mala intención, un artículo póstumo de Theodora Kallis lleno de sus obsesiones: el díctamo, los cretismos, una teoría sobre el libro VII de Safo. Lucía entiende que los poemas salen del círculo de Theodora. Busca quién la rodeaba y encuentra el nombre de Delia en los agradecimientos.
No la denuncia. La observa.
Entonces llega el séptimo poema. En él hay, transpuesto al griego eólico, un verso del libro olvidado de Lucía, el que escribió a los veinticinco años. Nadie más en el mundo podría reconocerlo. Lucía comprende tres cosas a la vez: Delia es la falsificadora, los poemas cambiaron de destinataria en algún momento de la serie, y la nueva destinataria es ella.
Relee los siete poemas buscando el punto de cambio y lo encuentra: a partir del cuarto la voz se dirige a una persona viva. Eso coincide con la fecha en que la contrataron.
Tercer acto
Lucía se enfrenta a Delia. Delia confiesa, pero se niega a decir lo único que importa: si los últimos poemas son una declaración de amor o una manera de asegurarse de que la experta que tenía que autentificarlos no pudiera traicionarla. Ella misma no está segura. Jura que el primer fragmento, el de Theodora, era auténtico, y que todo lo demás nació de ahí. No hay forma de comprobarlo.
Achterberg intuye que Lucía sabe algo y presiona: si el informe final no avala los poemas, Delia irá a prisión por el robo de papiro, y él se encargará de que así sea. Pell, por su parte, ha conseguido una muestra de tinta y anuncia unos resultados que nadie sabe si son concluyentes. Weller solo quiere su inauguración.
Delia desaparece y deja a Lucía un octavo papiro, que según ella es el último. Es el único escrito sabiendo que Lucía lo leería sabiendo quién lo escribió. Ya no habla como Safo, ni como Delia.
Final abierto
Es el día de la inauguración. Lucía lleva en el bolso dos informes, uno que autentifica los poemas y otro que los desenmascara, y el octavo papiro sin traducir. Lo ha leído en griego, pero no lo ha pasado a su lengua, porque traducirlo sería responderle.
Sube al estrado. La novela termina cuando abre la boca, o una página antes, en el hotel, de madrugada, cuando por primera vez en veinte años escribe un verso propio. No sabemos si es una traducción o una respuesta.
Quedan sin cerrar cuatro cosas: si el primer fragmento era de Safo, si Delia la amaba o la necesitaba, qué dice Lucía en el estrado y dónde está Delia.